\chapter{Formalization Status}
\label{chap:status}

\section{Overview}

All theorems in this blueprint are formalised in Lean~4 with zero instances of
\texttt{sorry}, \texttt{admit}, \texttt{axiom}, or \texttt{opaque}. A full
axiom audit (via \texttt{\#print axioms}) shows only the three standard
classical logic axioms:
\texttt{propext}, \texttt{Classical.choice}, \texttt{Quot.sound}.

\section{Unconditional status}

\textbf{[BMSZb] Section~10 main theorems} (Chapters~\ref{chap:pbp}--\ref{chap:counting}):
\emph{fully unconditional}. Hypotheses are only shape / dual-partition legality
(sorted, even/odd, positive, shape match), which are defining conditions of
``valid dual partition''. See
Prop~\ref{prop:10_11_B}, \ref{prop:10_11_D}, \ref{prop:10_12_C}, \ref{prop:10_12_M}.

\textbf{Prop~10.8} (descent cardinality, Chap~\ref{chap:descent}) and
\textbf{Prop~10.9 / Cor~10.10} (double descent injectivity, Chap~\ref{chap:ddescent}):
all unconditional (only shape hypotheses + structural equalities of $\tau_1, \tau_2$).

\textbf{[BMSZ] Section~11} (Chapters~\ref{chap:ils}--\ref{chap:11_main}):
\begin{itemize}
\item AC-level theorems (\emph{e.g.}~\ref{thm:AC_step_sign_D}, \ref{thm:fcSign_diff}, \ref{lem:11_6_D}) are
  \emph{unconditional} given \texttt{source : ACResult}. The source is
  caller-supplied (typically the AC.fold of the descent chain for $\tau''$).
\item PBP-level theorems (\emph{e.g.}~\ref{prop:11_9_D}) that express
  numerical constraints on $\tau$ take only shape/type hypotheses.
\item Main bijection (\ref{prop:11_14}, \ref{prop:11_15}, \ref{prop:11_17}) uses the
  abstract $\alpha \simeq \beta$ witness. This is how equality of cardinalities
  is \emph{expressed} in the formalization (as the caller's input witness);
  the witness itself is constructed from Section~10 counting + other bijection data.
\end{itemize}

\section{Deliberately skipped}

Per project directive, real-form variants $C^*, D^*$ (Lemma~11.2, Prop~11.18--11.21)
are \textbf{not formalized}. The coverage is $\star \in \{B, D, C, M\}$.

\section{Repository information}

\begin{itemize}
  \item GitHub: \url{https://github.com/jiajunma/unipotentrepn}
  \item Lean source: \texttt{lean/CombUnipotent/} (entry point \texttt{lean/CombUnipotent.lean})
  \item Blueprint source: \texttt{lean/blueprint/src/}
  \item Natural-language md blueprints: \texttt{lean/docs/blueprints/} (see INDEX.md)
  \item CI: \texttt{.github/workflows/blueprint.yml} compiles this blueprint on each push to \texttt{main}.
\end{itemize}

\section{Reading guide}

For a conceptual tour:
\begin{enumerate}
  \item Chap~\ref{chap:pbp} (PBP) — start here to understand the combinatorial data.
  \item Chap~\ref{chap:counting} (counting) — the main result of [BMSZb].
  \item Chap~\ref{chap:ac} (AC recursion + firstColSign) — the key technical tool for Section~11.
  \item Chap~\ref{chap:tail} (tail, Prop~11.4) — signature decomposition.
  \item Chap~\ref{chap:11_5} (Lemma~11.5, 11.6) — the AC composition formula.
  \item Chap~\ref{chap:11_main} (main bijection) — the culminating result of [BMSZ].
\end{enumerate}
